\documentclass[11pt,a4paper]{article}

%% PACHETE MATEMATICE
\usepackage{amsmath,amssymb,amsfonts}
\usepackage{amsthm}
\usepackage{mathpazo}
\usepackage{eulervm}
\usepackage{enumerate}
\usepackage{color}
\usepackage{graphicx}
\usepackage[all]{xy}
\usepackage{xcolor}
\usepackage{framed}
\usepackage[unicode]{hyperref}
\setcounter{MaxMatrixCols}{20}
\usepackage{multicol}
%\usepackage[romanian]{babel}


%% ASPECT
\linespread{1.05}
\usepackage[T1]{fontenc}
\usepackage[utf8x]{inputenc}
\usepackage[margin=2cm]{geometry}
% \usepackage[color]{showkeys}
\usepackage{anyfontsize}
\usepackage{titlesec}
\titleformat*{\section}{\large\bfseries}

\usepackage{fancyhdr}
\pagestyle{fancy}
\rhead{\small \color{gray} Notițe de Adrian Manea}
\lhead{\small \color{gray} IntProb-FIA, 2017---2018, L. Lemnete \& M. Olteanu}


\usepackage[autostyle = false, style = german]{csquotes}
\usepackage[romanian]{babel}
\newcommand\qq{\enquote}

%% COMENZILE MELE
\renewcommand*{\proofname}{Dem.:}
\newcommand\bb{\mathbb}
\newcommand\dr{\mathrm}
\newcommand\kal{\mathcal}
\newcommand\fk{\mathfrak}
\newcommand\op{\oplus}
\newcommand\ot{\otimes}
\newcommand\ds{\displaystyle}
\newcommand\id{\indent}
\newcommand\nid{\noindent}
\newcommand\seq{\subseteq}
\newcommand\sneq{\subsetneq}
\newcommand\speq{\supseteq}
\newcommand\spneq{\supsetneq}
\newcommand\rar{\rightarrow}
\newcommand\Rar{\Rightarrow}
\newcommand\xrar{\xrightarrow}
\newcommand\lar{\leftarrow}
\newcommand\Lar{\Leftarrow}
\newcommand\xlar{\xleftarrow}
\newcommand\lrar{\leftrightarrow}
\newcommand\Lrar{\Leftrightarrow}
\newcommand\ideal{\trianglelefteq}
\newcommand\ai{\text{ a.\^{i}. }}
\newcommand\sk{(\textit{Schi\c{t}\u{a}}) }
\newcommand\rhup{\rightharpoonup}
\newcommand\rhdn{\rightharpoondown}
\newcommand\lhup{\leftharpoonup}
\newcommand\lhdn{\leftharpoondown}
\newcommand\vid{\emptyset}
\newcommand\wh{\widehat}
\newcommand\ol{\overline}
\newcommand\NN{\mathbb{N}}
\newcommand\RR{\mathbb{R}}
\newcommand\ZZ{\mathbb{Z}}
\newcommand\QQ{\mathbb{Q}}
\newcommand\CC{\mathbb{C}}

%% STILURILE MELE
\newtheoremstyle{dotless}{}{}{\itshape}{}{\bfseries}{:}{ }{}

\theoremstyle{dotless}
\newtheorem{theorem}{Teorem\u{a}}[section]
\newtheorem{proposition}{Propozi\c{t}ie}[section]
\newtheorem{lemma}{Lem\u{a}}[section]
\newtheorem{corollary}{Corolar}[section]
\newtheorem{exercise}{Exerci\c{t}iu}

\newtheoremstyle{dotlessdr}{}{}{}{}{\bfseries}{:}{ }{}
\theoremstyle{dotlessdr}
\newtheorem{example}{Exemplu}[section]
\newtheorem{definition}{Defini\c{t}ie}[section]
\newtheorem{remark}{Observa\c{t}ie}[section]




\begin{document}

\begin{center}
  {\large\textbf{Seminar 4 \\ Integrale multiple}}
\end{center}

\vspace{1cm}

\section{Măsură și integrală Lebesgue}

\indent\indent Pentru integrarea în mai multe dimensiuni, trebuie să avem un concept
de \emph{măsură}. Dacă integrala Riemann permite calculul de lungimi, arii și volume
pentru o funcție de o variabilă reală, într-un spațiu multidimensional, este necesar
să adaptăm noțiunile respective, pentru a putea vorbi în continuare de \qq{cantități %
  măsurabile}. Acestea vor fi generalizări ale conceptelor de tip \emph{lungime, arie,
  volum}, în sensul că, în cazul particular al două sau trei dimensiuni, vor coincide
cu intuiția geometrică.

Fără a intra în detalii suplimentare, menționăm că măsura care se utilizează în acest
caz al problemelor din $\RR^n$ se numește \emph{măsura Lebesgue}. În particular,
\emph{elementul de lungime} față de măsura Lebesgue va fi notat cu $dl$ sau $dx$,
\emph{elementul de arie}, cu $dA = dxdy$, iar \emph{elementul de volum}, cu
$dV = dxdydz$ sau corespunzător unor eventuale schimbări de coordonate (e.g.\ în
coordonate polare, $dA = drd\theta$, în coordonate sferice, $dV = dr d\theta d\varphi$
etc.).

Față de această măsură Lebesgue, putem defini \emph{integrala Lebesgue} din $\RR^k$,
a unei funcții reale $f : A \seq \RR^k \to \RR$, notată, în general, prin:
\[
  \int_A f(x_1, \dots, x_k) dx_1 \dots dx_k.
\]

În cazurile particulare $k = 1, 2, 3$, vom folosi notațiile uzuale:
\[
  \int_A f(x) dx, \quad \iint_A f(x, y) dxdy, \quad \iiint_A f(x, y, z) dxdydz.
\]

Următorul rezultat ne arată că putem inversa ordinea de integrare:
\begin{theorem}[Fubini]\label{thm:fubini}
  Fie $(x, y) \in \RR^{k + p}$ un punct oarecare, notăm măsura Lebesgue din $\RR^k$ cu
  $dx$, unde $x \in \RR^k$, măsura Lebesgue din $\RR^p$ cu $dy, y \in \RR^p$, iar
  măsura Lebesgue din $\RR^{k + p}$ cu $dxdy$.

  Fie $f : \RR^{k + p} \to \RR$ o funcție integrabilă Lebesgue. Atunci are loc:
  \[
    \int_{\RR^k} \Bigg( \int_{\RR^p} f(x, y) dy \Bigg) dx = %
    \iint_{\RR^{k + p}} f(x, y) dxdy = %
    \int_{\RR^p} \Bigg( \int_{\RR^k} f(x, y) dx \Bigg) dy.
  \]
\end{theorem}

Cazurile particulare pe care le vom folosi în exerciții sînt:
\begin{enumerate}[(a)]
\item \emph{Aria} unei mulțimi compacte $D \seq \RR^2$ se calculează prin:
  \[
    \kal{A}(D) = \iint_D dxdy;
  \]
\item \emph{Volumul} unei mulțimi compacte $\Omega \seq \RR^3$ se calculează prin:
  \[
    \kal{V}(\Omega) = \iiint_\Omega dxdydz.
  \]
\end{enumerate}

De asemenea, pentru calculul integralelor multiple, putem adapta formula de
schimbare de variabile:

\begin{theorem}[Schimbare de variabile]\label{thm:sch-var}
  Fie $A \seq \RR^n$ o mulțime deschisă și $\varphi$ un difeomorfism\footnote{%
    aplicație continuă și bijectivă, cu inversa continuă} al lui $A$.

  Pentru orice funcție continuă $f : \varphi(A) \to \RR$ are loc:
  \[
    \int_{\varphi(A)} f(x) dx = \int_A (f \circ \varphi)(y) \cdot |J_\varphi(y)| dy,
  \]
  unde $J_\varphi$ este jacobianul difeomorfismului (schimbării de variabile) $\varphi$.
\end{theorem}

%%%%%%%%%%%%%%%%%%%%%%%%%%%%%%%%%%%%%%%%%%%%%%%%%%%%%%%%%%%%%%%%%%%%%%

\section{Exerciții}

1. Să se calculeze următoarele integrale duble:
\begin{enumerate}[(a)]
\item $\ds\iint_D xy^2 dxdy$, unde $D = [0, 1] \times [2, 3]$;
\item $\ds\iint_D xy dxdy$, unde $D = \{(x, y) \in \RR^2 \mid y \in [0,1], y^2 \leq x \leq y \}$;
\item $\ds\iint_D ydxdy$, unde $D = \{(x, y) \in \RR^2 \mid (x - 2)^2 + y^2 \leq 1 \}$.
\end{enumerate}

\emph{Indicație:} Integralele pot fi gîndite ca integrale iterate. În particular, pentru
domeniul de la punctul (c), avem $-\sqrt{(x - 2)^2 - 1} \leq y \leq \sqrt{(x - 2)^2 - 1}$,
iar $x \in [1, 3]$ (altfel, inegalitatea nu are soluții).

\vspace{1cm}

2. Să se calculeze integralele duble:
\begin{enumerate}[(a)]
\item $\ds\iint_D (x + 3y) dxdy$, unde $D$ este mulțimea plană mărginită de curbele de
  ecuații:
  \[
    y = x^2 + 1, \quad y = -x^2, \quad x = -1, \quad x = 3;
  \]
\item $\ds\iint_D e^{|x + y|} dxdy$, unde $D$ este mulțimea plană mărginită de curbele
  de ecuații:
  \[
    x + y = 3, \quad x + y = -3, \quad y = 0, \quad y = 3;
  \]
\item $\ds\iint_D xdxdy$, unde $D$ este mulțimea plană mărginită de $x^2 + y^2 = 9, x \geq 0$.
\end{enumerate}

\emph{Indicație:} În toate cazurile, este util să schițăm grafic mulțimea, pentru a
vedea care sînt capetele domeniilor pentru $x$, respectiv $y$. În plus, în cazul (b),
avem de explicitat modulul, deci studiem cazurile cînd $x + y \geq 0$ și $x + y < 0$
pe domeniile de definiție.

\vspace{1cm}

3. Folosind trecerea la coordonate polare, să se calculeze integralele:
\begin{enumerate}[(a)]
\item $\ds\iint_D e^{x^2 + y^2} dxdy$, unde $D: x^2 + y^2 \leq 1$;
\item $\ds\iint_D (1 + \sqrt{x^2 + y^2}) dxdy$, unde $D : x^2 + y^2 - y \leq 0, x \geq 0$;
\item $\ds\iint_D \ln(1 + x^2 + y^2) dxdy$, unde $D$ este mărginit de curbele de ecuații:
  \[
    x^2 + y^2 = e^2, \quad y = x\sqrt{3}, \quad x = y \sqrt{3}, \quad x \geq 0.
  \]
\end{enumerate}

\vspace{1cm}

4. Calculați integralele duble, găsind explicit domeniul de definiție (prin intersecția
curbelor care îl definesc):
\begin{enumerate}[(a)]
\item $\ds\iint_D xydxdy$, cu $D: xy = 1, \quad x + y = \frac{5}{2}$;
\item $\ds\iint_D \frac{1}{\sqrt{x}} dxdy$, cu $D: y^2 \leq 8x, y \leq 2x, y + 4x \leq 24$;
\item $\ds\iint_D (1 - y) dxdy$, cu $D: x^2 + (y-1)^2 \leq 1, y \leq x^2, x \geq 0$.
\end{enumerate}


\end{document}